\section{Semiconductor Photodetector Theory}

![[Pasted image 20260814160058.png]] In a \textbf{Schottky photodiode}, light with
energy $h\nu > E_{\mathrm{g}}$, i.e. above the \emph{band gap energy}
$E_{\mathrm{g}}$, excites an \emph{electron-hole pair}, such that the excited electron
resides in the conduction band while the excited hole exists in the valence band.
The excitation may occur inside the \emph{space-charge region} (SCR) or in the neutral
$n$-region. Inside the SCR, the built-in electric field separates both charge carriers
and creates a \emph{drift current}, such that only holes are transported through the
metal contact, while electrons move towards the neutral $n$-region. Charge carriers
generated in the $n$-region may recombine. As there is no electric field present
in this region, currents only establish through diffusion.

The Schottky diode is one-dimensional along the $z$ direction, with cross
section $A$. It features three distinct regions as:

\[
	\underbrace{ z < 0 }_{ \text{metal contact} }\quad | \underbrace{ \quad 0 \leq z
	\leq z_{0} \quad }_{ \text{space-charge region (SCR)} }| \quad \underbrace{ z >
	z_{0} }_{ \text{quasi-neutral } n\text{-bulk} }
\]
Light enters from $z<0$ and propagates in the $+z$ direction. The space charge
region directly begins at the semiconductor surface.

\subsection{Physical Modelling}

The goal is to find an expression for the \emph{photocurrent}, depending on the photon
generation rate. To this end, a differential equation is created using the Electron-Hole
Continuity Equation, together with the definition of the Semiconductor Current
Density|Current Densities. The central idea is to utilize the Conservation of Total
Current to probe the current specifically at the boundary between SCR and neutral
region ($z=z_{0}$) to obtain the current anywhere through the diode.

\subsubsection{Boundary Conditions}

As the electric field points towards the metal contact, electrons generated in the
depletion layer are transported in $+z$ direction. Only holes leave the device
through this contact.

Let $p_{0}$ be the equilibrium hole concentration (without any optical excitation).
The electric field is assumed to extract all excess holes that arrive at $z_{0}$,
acting as a perfect sink. Far away from the junction, the bulk is undisturbed, as
effectively all photons are absorbed beforehand.

\subsubsection{Continuity Equation in Depletion Layer}

The goal is to evaluate the electron current density at $z_{0}$. The electron Electron-Hole
Continuity Equation|Continuity Equation reads:

\[
	\partial_{t}n(z) = G_{\mathrm{n}}(z) - U_{\mathrm{n}}(z) + \frac{1}{e}\partial_{z}
	j_{\mathrm{n}}(z)\\
\]
As the electric field is strong and transports electrons away rapidly,
recombination in this region can be neglected. Additionally assuming \emph{steady
state} reduces the equation to:

\[
	\begin{align}
		\partial_{z}j_{\mathrm{n}}              & =-eG(z)                                \\
		j_{\mathrm{n}}(z_{0})-j_{\mathrm{n}}(0) & =-e \int_{0}^{z_{0}}G(z) \,\mathrm{d}z
	\end{align}
\]
Here, $G$ denotes the generation rate of electron-hole pairs, which is specified
in the last section. Applying the boundary condition $j_{\mathrm{n}}(0)=0$ yields:

\[
	j_{\mathrm{n}}(z_{0})=-e \int_{0}^{z_{0}}G(z) \, \mathrm{d}z
\]
\subsubsection{Continuity Equation in Bulk Layer}

The goal is to evaluate the hole current density at $z_{0}$. Define $\Delta p(z)$
as the \emph{excess charge carrier density}, such that $p(z) = p_{0}+ \Delta p(z)$
with $p_{0}$ being the constant equilibrium charge carrier density. Substituting
the Semiconductor Current Density into the hole Electron-Hole Continuity
Equation|Continuity Equation results in the transport equation:

\[
	\begin{align}
		\partial^{2}_{z}\Delta p(z)-\frac{eE}{kT}\partial_{z}\Delta p(z)=\frac{U(z)-G(z)}{D_{\mathrm{p}}} \\
	\end{align}
\]
At \emph{low injection} the net recombination is linear in $\Delta p$, with
proportionality constant $\tau_{\mathrm{p}}$:

\[
	U(z)=\frac{\Delta p}{\tau_{\mathrm{p}}}
\]
This leads to a linear, inhomogeneous second order differential equation:

\[
	\partial_{z}^{2}\Delta p(z)-\frac{eE}{kT}\partial_{z}\Delta p(z)-\frac{\Delta
	p(z)}{D_{\mathrm{p}}\tau_{\mathrm{p}}}=-\frac{G(z)}{D_{\mathrm{p}}}
\]
As the electric field is present only in the Depletion Layer, the second term
vanishes:

\[
	\partial_{z}^{2}\Delta p-\frac{\Delta p}{D_{\mathrm{p}}\tau}=-\frac{G(z)}{D_{\mathrm{p}}}
\]
Define $\lambda=1 / \sqrt{ D_{\mathrm{p}}\tau_{\mathrm{p}}}$ and use $'$ to
indicate derivatives. Multiplication of the left hand side with $\exp(-\lambda(z-
z_{0}))$ is equivalent to differentiation:

\[
	\begin{align}
		 & \frac{\mathrm{d}}{\mathrm{d}z}[\exp(-\lambda(z-z_{0})) (\Delta p'+\lambda\Delta p)]                                \\
		 & =-\lambda \exp(-\lambda(z-z_{0}))(\Delta p'+\lambda\Delta p)+\exp(-\lambda(z-z_{0})) (\Delta p''+\lambda\Delta p') \\
		 & =\exp(-\lambda(z-z_{0})) (\Delta p'' -\lambda^{2}\Delta p) =-\frac{G(z)}{D_{\mathrm{p}}}\exp(-\lambda(z-z_{0}))
	\end{align}
\]
Integration from $z_{0}$ to $\infty$ yields:

\[
	[\exp(-\lambda(z-z_{0})) (\Delta p'+\lambda\Delta p)]^{\infty}_{z_{0}}=-\frac{1}{D_{\mathrm{p}}}
	\int_{z_{0}}^{\infty}G(z) \exp(-\lambda(z-z_{0})) \, \mathrm{d}z
\]
For $z \to \infty$, the exponential term vanishes, while at $z=z_{0}$, it
reduces to 1 and $\Delta p$ vanishes. The equation simplifies to:

\[
	D_{\mathrm{p}}\Delta p'(z_{0})=\int_{z_{0}}^{\infty}G(z) \exp(-\lambda(z-z_{0})
	) \, \mathrm{d}z
\]
As there is no electric field, current only originates from diffusion:

\[
	j_{\mathrm{p}}=-eD_{\mathrm{p}}\partial_{z}p \\=-e \int_{z_{0}}^{\infty}G(z) \exp
	(-\lambda(z-z_{0})) \, \mathrm{d}z
\]
Note that the homogeneous solution of the differential equation
$\partial_{z}^{2}\Delta p-\frac{\Delta p}{D_{\mathrm{p}}\tau}=0$ together with
the boundary condition $\Delta p(z \to \infty)$ is given as:

\[
	\Delta p=\Delta p_{0}\exp(-\lambda(z-z_{0}))
\]
Therefore, it is justified to write:

\[
	j_{\mathrm{p}}=-e \int_{z_{0}}^{\infty}G(z) \frac{\Delta p}{\Delta p_{0}}\, \mathrm{d}
	z
\]
\subsubsection{Current Characteristics}

The total current density is defined as: $j_{\mathrm{tot}}=j_{\mathrm{n}}(z_{0})+
j_{\mathrm{p}}(z_{0})$, substituting results from previous chapters yields:

\[
	\begin{align}
		j_{\mathrm{tot}} & =-e \int_{0}^{z_{0}}G(z)\, \mathrm{d}z - e\int_{z_{0}}^{\infty}G(z) \exp(-\lambda(z-z_{0})) \, \mathrm{d}z                                                           \\
		                 & =-e \int_{0}^{\infty}G(z) j(z)\, \mathrm{d}z \qquad j(z)=\left\{ \begin{matrix}1&z \in(0, z_{0}] \\ \exp(-\lambda(z-z_{0}))&z \in(z_{0}, \infty)\end{matrix} \right.
	\end{align}
\]
The quantity $j(z)$ is called the \emph{point source current} induced by the point
source at position $z$. Multiplication with the cross section $A$ results in an
expression for the current:

\[
	I_{\mathrm{ph}}=Aj_{\mathrm{tot}}=-eA \int_{0}^{\infty}G(z)j(z) \, \mathrm{d}z
\]
Note that the current originating from the SCR at $z_{0}$ is a pure electron
current, while the hole current from the neutral current originates from holes only.

\subsubsection{Definition of the Generation Rate}

The generation rate $G$ of electron-hole pairs is proportional to the photon
absorption rate $U_{\mathrm{ph}}$, with the \emph{internal quantum efficiency}
$\eta$ as proportionality constant $G=\eta U_{\mathrm{ph}}$. Due to the Photon Continuity
Equation $U_{\mathrm{ph}}=-\frac{\mathrm{d} \Phi}{\mathrm{d}z}$, where $\Phi$ is
the \emph{photon flux}. An expression for the flux is given by the Lambert-Beer
Law, which leads to:

\[
	G(z)=\eta \Phi_{0}\alpha(h\nu, z)\exp\left( -\int_{0}^{z}\alpha(h \nu, z') \, \mathrm{d}
	z' \right)
\]